%% ============================================================
%% Episode 11 - Farsi Podcast Script (Translation)
%% Source: specialized_advisor_report_complete/markdown/11_missing_sections_and_additions.md
%% Series: Advisor Progress Report - Deep Dive
%% Compiler: XeLaTeX ONLY
%% Font: B Nazanin
%% CRITICAL: xepersian must always be loaded LAST
%% ============================================================

\documentclass[12pt, a4paper]{article}

%% ============================================================
%% COMPATIBILITY SHIMS
%% ============================================================
\makeatletter
\providecommand\IfClassLoadedT[2]{\@ifclassloaded{#1}{#2}{}}
\providecommand\IfClassLoadedF[2]{\@ifclassloaded{#1}{}{#2}}
\providecommand\IfPackageLoadedT[2]{\@ifpackageloaded{#1}{#2}{}}
\providecommand\IfPackageLoadedF[2]{\@ifpackageloaded{#1}{}{#2}}
\providecommand\IfFileLoadedTF[3]{\@ifpackageloaded{#1}{#2}{#3}}
\providecommand\IfFileLoadedT[2]{\@ifpackageloaded{#1}{#2}{}}
\providecommand\IfFileLoadedF[2]{\@ifpackageloaded{#1}{}{#2}}
\makeatother
\usepackage{expl3}
\ExplSyntaxOn
\cs_if_exist:NF \prop_gput_if_not_in:Nnn
  {
    \cs_new:Npn \prop_gput_if_not_in:Nnn #1#2#3
      { \prop_if_in:NnF #1 {#2} { \prop_gput:Nnn #1 {#2} {#3} } }
  }
\ExplSyntaxOff
%% ============================================================

\usepackage[
  a4paper,
  top=2.5cm, bottom=2.5cm,
  left=2.8cm, right=2.8cm,
  headheight=25pt
]{geometry}

\usepackage{xcolor}
\definecolor{seriesblue}{RGB}{25, 80, 150}
\definecolor{audionote}{RGB}{130, 50, 15}
\definecolor{sectiongray}{RGB}{75, 75, 75}
\definecolor{audiotan}{RGB}{252, 243, 233}

\usepackage{amsmath}
\usepackage{amssymb}
\usepackage{enumitem}
\usepackage{setspace}
\setstretch{1.8}
\usepackage{mdframed}
\usepackage{titlesec}

%% ============================================================
%% xepersian MUST be loaded LAST
%% ============================================================
\usepackage{xepersian}
\settextfont[Scale=1.05]{B Nazanin}
\setlatintextfont{Times New Roman}

\newcommand{\dash}{\lr{---}}

%% ============================================================
%% Page style -- savebox approach
%% ============================================================
\newsavebox{\hdrLeft}
\newsavebox{\hdrRight}
\AtBeginDocument{%
  \savebox{\hdrLeft}{\normalfont\small\color{sectiongray}%
    گزارش پيشرفت مشاور / بررسی عمیق}%
  \savebox{\hdrRight}{\normalfont\small\color{sectiongray}%
    قسمت يازدهم: شکاف‌های تکميل}%
}
\makeatletter
\def\ps@podcast{%
  \def\@oddhead{\usebox{\hdrLeft}\hfill\usebox{\hdrRight}}%
  \let\@evenhead\@oddhead
  \def\@oddfoot{\normalfont\hfil\thepage\hfil\relax}%
  \let\@evenfoot\@oddfoot
  \let\@mkboth\markboth
}
\makeatother
\pagestyle{podcast}

%% ============================================================
%% Section formatting
%% ============================================================
\titleformat{\section}
  {\large\bfseries\color{seriesblue}}
  {}
  {0pt}
  {}
  [\vspace{-2pt}{\color{seriesblue}\rule{\linewidth}{0.7pt}}\vspace{2pt}]

\titleformat{\subsection}
  {\normalsize\bfseries\color{sectiongray}}
  {}
  {0pt}
  {}

\setcounter{secnumdepth}{0}

\setlist[itemize]{
  itemsep=4pt,
  topsep=5pt,
  parsep=0pt,
  label=$\bullet$
}

%% ============================================================
\begin{document}
%% ============================================================

\thispagestyle{empty}

\vspace*{1.2cm}

{\centering
{\color{seriesblue}\fontsize{24}{30}\selectfont\bfseries قسمت يازدهم}

\vspace{0.5cm}
{\fontsize{15}{20}\selectfont\bfseries شکاف‌های تکميل پايان‌نامه: ده درصد باقی‌مانده}

\vspace{0.9cm}
{\color{seriesblue}\rule{0.55\linewidth}{1pt}}
\vspace{0.7cm}

{\normalsize
\begin{tabular}{rl}
  \textbf{مجموعه:}     & گزارش پيشرفت مشاور \dash{} بررسی عمیق \\[5pt]
  \textbf{مدت:}        & ۸ تا ۱۰ دقیقه \\[5pt]
  \textbf{راوی:}       & مجری تنها \\[5pt]
  \textbf{منابع:}      & فهرست بررسی حسابرسی گزارش مشاور \\
\end{tabular}
}

\vspace{0.9cm}
{\color{seriesblue}\rule{0.55\linewidth}{0.4pt}}
\par}

\vspace{1.2cm}

%% ============================================================
\section{افتتاحیه: ۹۰ از ۱۰۰}

امتياز تکميل پايان‌نامه مستند در گزارش ۹۰ از ۱۰۰ است. يعنی ۹۰ درصد کار انجام شده و ۱۰ درصد باقيمانده شناسايی، به مراحل تقسيم، و تاريخ هدف داده شده.

اين قسمت مشخص می‌کند آن ۱۰ درصد باقيمانده در عمل چگونه است \dash{} شش شکاف خاص که بايد قبل از تکميل پايان‌نامه پر شوند. درک اينها نه فقط برای برنامه‌ريزی تکميل بلکه برای گفتگوهای مشاور مهم است: شکاف‌ها به صداقت در گزارش مستند شده‌اند، و يک مشاور ممکن است آن‌ها را مستقيماً بررسی کند.

%% ============================================================
\section{شکاف اول: جدول مقايسه با ادبيات}

گزارش نتايج عملکرد برای چهار کنترل‌گر در يک سيستم آونگ معکوس مضاعف تأييد می‌کند. اما آن نتايج را با کارهای منتشرشده از محققان ديگر در سيستم‌های مشابه يا يکسان مقايسه نمی‌کند.

يک جدول معيار ادبيات نشان می‌داد: اينجاست که بهترين کنترل‌گرهای منتشرشده به آن رسيدند، اينجاست که اين کار به آن می‌رسد، و اينجاست که اين کار در برابر وضعيت هنر قرار می‌گيرد.

بدون اين مقايسه، يک سوال کليدی بی‌پاسخ می‌ماند: اين نتايج خوب هستند؟ آيا ۱.۸۲ ثانيه زمان يکنواختی برای \lr{STA} با ادبيات منتشرشده رقابتی است، يا ناقابل توجه؟

چالش: يافتن مقالاتی که پيکربندی‌های \lr{DIP} يکسان يا قابل مقايسه را مطالعه کرده‌اند \dash{} همان تعداد لينک‌ها، همان رژيم پارامتر، همان معيارهای عملکرد. مقايسه منصفانه نيازمند تطابق دقيق شرايط آزمايش است. حداقل نسخه قابل دفاع: سه تا پنج مقاله مرتبط را نقل کنيد، توجه کنيد زمان‌های يکنواختی و معيارهای لرزش گزارش‌شده‌شان چيست، و تأييد کنيد کجا مقايسه مستقيم ممکن است و کجا نيست.

%% ============================================================
\section{شکاف دوم: آناليز شبيه‌سازی به سخت‌افزار}

هر نتيجه در اين پروژه از شبيه‌سازی می‌آيد. گزارش درباره اين صريح است. اما وقتی اين کنترل‌گر روی سخت‌افزار واقعی اجرا می‌شود چه می‌شود؟

سخت‌افزار واقعی اثراتی معرفی می‌کند که در شبيه‌سازی غايب هستند: ديناميک اشباع محرک \dash{} ثابت زمانی بين فرمان دادن به يک نيرو و توليد آن توسط موتور. کوانتيزاسيون انکودر \dash{} اندازه‌گيری‌های زاويه ديجيتال هستند، نه پيوسته. تأخيرهای ارتباطی بين حسگر و کنترل‌گر. اصطکاک کولمب و چسبندگی که با مدل متفاوتند. تشديدهای ساختاری در لينک‌های فيزيکی.

آناليز شبيه‌سازی به سخت‌افزار نيازی به ساخت واقعی سخت‌افزار ندارد. نياز دارد به شناسايی سيستماتيک فرض‌ها در مدل شبيه‌سازی که بيشترين احتمال شکستن روی سخت‌افزار را دارند، و سپس يا تحليل اينکه آيا کنترل‌گرها به هر اثر مقاوم هستند، يا هرکدام را به عنوان محدوديت شناخته‌شده با تخفيف پيشنهادی علامت‌گذاری کنيد.

شش اثر که بايد پرداخته شوند: ديناميک‌ها و اشباع محرک، کوانتيزاسيون انکودر، تأخيرها و نويز، اصطکاک و چسبندگی مدل‌نشده، انعطاف ساختاری و تشديد، و لرزش زمان‌بندی بلادرنگ.

%% ============================================================
\section{شکاف سوم: هزينه محاسباتی بهينه‌سازی ازدحام ذرات}

گزارش پيکربندی \lr{PSO} را مستند می‌کند \dash{} ۴۰ ذره، ۲۰۰ تکرار، کنترل‌گرها دقيقه‌ها می‌گيرند. اما زمان‌های ساعت ديوار خاصی به ازای هر کنترل‌گر نمی‌دهد.

اين به دو دليل مهم است. اول، تکرارپذيری: اگر کسی بخواهد نتايج \lr{PSO} را تکرار کند، دانستن زمان محاسبه مورد انتظار به او کمک می‌کند برنامه‌ريزی کند. دوم، مقايسه: انتخاب مدل ساده‌شده در برابر کامل برای \lr{PSO} توسط سرعت انگيزه گرفت \dash{} اما چقدر سريع‌تر؟

آيتم‌های خاص مورد نياز: زمان ساعت ديوار \lr{PSO} به ازای هر کنترل‌گر روی دستگاه مرجع، مشخصات \lr{CPU} و حافظه، و مقايسه زمان شبيه‌سازی مدل ساده‌شده در برابر کامل برای توجيه ادعای سرعت.

اين صرفاً يک کار اندازه‌گيری و مستندسازی است، نه يک آزمايش جديد. اجراهای \lr{PSO} قبلاً انجام شدند. زمان‌ها فقط بايد ثبت و گزارش شوند.

%% ============================================================
\section{شکاف چهارم: تکرارپذيری مطالعه \lr{QW-2}}

مطالعه \lr{QW-2} در گزارش به عنوان تأييد تک-اجرايی پايه ظاهر می‌شود. فاصله‌های اطمينان برای هيبريد تطبيقی ابر-پيچشی \dash{} اعدادی که در جدول \lr{CI} دوم ظاهر می‌شوند \dash{} از اين مطالعه می‌آيند.

مشکل: بردار شرايط اوليه دقيق و تنظيمات سناريو برای \lr{QW-2} در متن اصلی بيان نشده‌اند. يک خواننده بدون دانستن وضعيت شروع نمی‌تواند نتيجه \lr{QW-2} را تکرار کند.

اين يک شکاف مستندسازی ساده است. افزودن سه خط به متن اصلی \dash{} بردار حالت اوليه استفاده‌شده در \lr{QW-2}، مدت شبيه‌سازی، و سطح نويز يا اغتشاش در صورت وجود \dash{} آن را کاملاً می‌بندد.

%% ============================================================
\section{شکاف پنجم: يکپارچه‌سازی ارجاع شکل‌ها}

گزارش ليستی از شکل‌های کامل‌شده دارد: نمودار معماری سيستم، شماتيک حلقه کنترل، پورتره فاز لياپانوف، سطح سه‌بعدی همگرايی \lr{PSO}. همه شکل‌ها وجود دارند و استاندارد کيفی ۳۰۰ \lr{DPI} را رعايت می‌کنند.

شکاف: نه همه شکل‌ها با عبارت‌های ارجاع صريح در متن ارجاع داده شده‌اند. قرارداد نوشتاری دانشگاهی مستلزم است هر شکلی حداقل يک ارجاع در متن داشته باشد \dash{} «همان‌طور که در شکل ۳ نشان داده شده» يا «برای پورتره فاز شکل ۵ را ببينيد.» اگر يک شکل ارجاع نداشته باشد، در سند شناور به نظر می‌رسد بدون ارتباط با روايت.

اين يک کار ويرايشی خط است، نه يک کار فنی. هر شکل را مرور کنيد، جايی در متن که بيشتر مرتبط است را پيدا کنيد، و جمله ارجاع را اضافه کنيد.

%% ============================================================
\section{شکاف ششم: بخش کار آينده اختصاصی}

گزارش در حال حاضر محدوديت‌ها و موارد باز را به صورت درون‌خطی مستند می‌کند \dash{} پراکنده در بخش‌های جداگانه به عنوان «محدوديت‌های شناخته‌شده»، «موارد باز»، و «افزودنی‌های مورد نياز». اين برای يادداشت پيشرفت مفيد است اما نه برای يک پايان‌نامه.

يک پايان‌نامه به يک بخش کار آينده اختصاصی نياز دارد که: همه موارد باز را در يک جا جمع کند، آن‌ها را بر اساس اهميت و امکان‌پذيری اولويت‌بندی کند، هر مورد را به مشارکت مربوطه پايان‌نامه که گسترش می‌دهد پيوند بزند، و بين موارد ضروری برای ادعاهای فعلی و مواردی که کار را گسترش می‌دهند تمييز قائل شود.

%% ============================================================

\begin{mdframed}[
  backgroundcolor=audiotan,
  linecolor=audionote!70,
  linewidth=1pt,
  innerleftmargin=12pt,
  innerrightmargin=12pt,
  innertopmargin=8pt,
  innerbottommargin=8pt
]
\textbf{\color{audionote}[يادداشت صوتی:]}
زمان‌بندی تکميل: مرحله ۳ \dash{} گسترش کتاب‌نامه به ۱۵ مرجع يا بيشتر، هدف ۱۰ ژانويه. مرحله ۴ \dash{} بازخوانی دقيق، هدف ۱۳ ژانويه. مرحله ۵ \dash{} بررسی تضمين کيفيت نهايی يا تأييد همه ارجاع‌های شکل، متقاطع-مراجع، و برچسب‌های معادله، هدف ۱۴ ژانويه. مرحله ۶ \dash{} بازبينی مشاور، هدف ۲۴ ژانويه. ۹۰ درصد کامل نزديک اتمام نيست. موارد خاص شناخته‌شده با تاريخ برای هر آيتم است. اين يک موضع قابل دفاع است.
\end{mdframed}

\vspace{0.6cm}

%% ============================================================
\section{نتيجه‌گيری}

نود درصد کامل تمام نيست. شش شکاف \dash{} جدول مقايسه ادبيات، آناليز شبيه‌سازی به سخت‌افزار، هزينه‌های محاسباتی \lr{PSO}، تکرارپذيری \lr{QW-2}، ارجاع‌های شکل، و بخش کار آينده \dash{} خاص و قابل اقدام هستند. هيچ‌کدام به آزمايش‌های جديد نياز ندارند. به آناليز، مستندسازی، و جستجوی ادبيات نياز دارند. سه تا چهار جلسه کاری متمرکز بايد همه را ببندد.

قسمت آخر: تمرين کامل دفاع از پايان‌نامه \dash{} هفت محتمل‌ترين سوال و چگونگی پاسخ به هر کدام.

%% ============================================================
\vspace{1.5cm}
\begin{center}
{\color{seriesblue}\rule{0.45\linewidth}{0.4pt}}\\[8pt]
{\small\color{sectiongray}
منابع: تحليل شکاف بر اساس فهرست بررسی حسابرسی \lr{advisor\_progress\_report.tex}
}
\end{center}

\end{document}
